%HOMEWORK template for M 325K
\documentclass{article}
\headheight=7pt         \topmargin=14pt
\textheight=574pt       \textwidth=445pt
\oddsidemargin=18pt     \evensidemargin=18pt

%Begin package import

\usepackage{amsmath, amssymb, amsthm}
\usepackage{mathtools}
\usepackage{pdfpages}
\usepackage{tikz-cd}

%End package import
%Begin custom commands

\newcommand*{\T}{\mathcal{T}}
\newcommand*{\HC}{\mathcal{H}}
\newcommand*{\B}{\mathcal{B}}
\newcommand*{\U}{\mathcal{U}}
\newcommand*{\cl}{\mathrm{cl}}
\newcommand*{\subeq}{\subseteq}
\newcommand*{\empset}{\emptyset}
\newcommand{\C}{\mathbb{C}}
\newcommand{\R}{\mathbb{R}}
\newcommand{\Q}{\mathbb{Q}}
\newcommand{\Z}{\mathbb{Z}}
\newcommand{\N}{\mathbb{N}}
\newcommand{\F}{\mathbb{F}}
\newcommand{\abs}[1]{\left\lvert #1\right\rvert}
\newcommand{\RM}[1]{\mathrm{#1}}
\newcommand{\BF}[1]{\textbf{#1}}
\newcommand{\e}{\epsilon}
\newcommand{\ceil}[1]{\lceil{#1}\rceil}
\newcommand{\floor}[1]{\lfloor{#1}\rfloor}
\newcommand{\rarr}[0]{\rightarrow}
\newcommand{\IM}[1]{\text{Im}(#1)}
\newcommand{\Hom}[0]{\text{Hom}}
\newcommand{\Pow}[1]{\text{Pow}(#1)}
\newcommand{\angles}[1]{\langle #1 \rangle}

%End custom commands
%Begin custom theorems

\newtheorem{innercustomgeneric}{\customgenericname}
\providecommand{\customgenericname}{}
\newcommand{\newcustomtheorem}[2]{%
\newenvironment{#1}[1]
{%
\renewcommand\customgenericname{#2}%
\renewcommand\theinnercustomgeneric{##1}%
\innercustomgeneric%
}
{\endinnercustomgeneric}
}

\newcustomtheorem{THRM}{Theorem}
\newcustomtheorem{LEM}{Lemma}
\newcustomtheorem{EX}{Exercise}
\newcustomtheorem{COR}{Corollary}
\newcustomtheorem{PROB}{Problem}
\newcustomtheorem{claim}{Claim}

\newenvironment{solution}
{\renewcommand\qedsymbol{$\blacksquare$}\begin{proof}[Solution]}
{\end{proof}}

\newenvironment{definition}
{\renewcommand\qedsymbol{$\blacksquare$}\begin{proof}[Definition]}
{\end{proof}}

%End custom theorems
\begin{document}

\title{Exercise 3}
\author{Arham Lodha}%put your name here
\date{March 28, 2025}%enter the due date here
\maketitle

\begin{proof}
    \begin{enumerate}
        \item $\implies$: Suppose $\forall g \in S_{k}(1)$, $a_{g}(n) = O(n^{(k - 1) / 2 + \epsilon})$. For Hecke Eigenforms $f \in H_{k} \subset S_k(1), a_f(n) = \lambda_f(n)$. By b, $\lambda_f(n) \leq d(n)n^{(k - 1)/2 + \epsilon}$. 
        \item $\impliedby:$ $H_k$ is a eigenbasis of $S_k(1)$ a finite vectorspace. Thus $\left\{ f_1, \dots, f_{\dim S_k} \right\} = H_k$ and $\forall g \in S_{k}(1)$
        
        \[g = c_{1} f_1 + \hdots + c_{\dim S_k} f_{\dim S_k}\]

        Thus $a_{g}(n) = \sum_{j = 1}^{\dim S_k} c_j a_{f_j}(n) = \sum_{j = 1}^{\dim S_k} c_j \lambda_{f_j}(n) \leq d(n)n^{(k - 1)/2} \sum_{j = 1}^{\dim S_k} c_j$ 
    \end{enumerate}
     
\end{proof}

\end{document}